\chapter{Bonos}\label{cap:bonos}

% -------------------------------------------------------------------------------------------------
\section{Propósito}\label{sec:cap06-proposito}

El producto \texttt{bond} de QF Desk es un bono de tasa fija con cupones iguales y amortización del principal
al vencimiento. El código calcula su precio sucio como suma de flujos descontados con una curva de la
aplicación (capítulo~\ref{cap:curvas}), y, para el bono envejecido (valuado después de su fecha de pacto), el
cupón corrido y el precio limpio. Además ofrece el rendimiento, la duración y la convexidad. No hay bonos
flotantes, amortizables ni con opcionalidad. El valor de un bono envejecido que entra al portafolio y al
P\&L se trata también en los capítulos~\ref{cap:ciclo-vida} (liquidación de cupones) y~\ref{cap:riesgo}
(riesgo).

% -------------------------------------------------------------------------------------------------
\section{Notación}\label{sec:cap06-notacion}

Los símbolos comunes están en el apéndice~\ref{ap:simbolos}; aquí $\Nom$ es el valor nominal (\texttt{face} en
el código), $\K$ la tasa de cupón anual (decimal) y $\DF{c}{T}$ el factor de descuento de la curva $c$ de
descuento del bono. Los símbolos propios son:
\begin{itemize}
  \item $f$: frecuencia (cupones por año); $m$: número de periodos (cupones) del bono; $t_i$: tiempo del pago
        $i$ en años desde la fecha de valuación.
  \item $\tau=\texttt{periodFraction}(f,\mathrm{dcc})$: fracción de devengo de un periodo uniforme (años);
        $C=\Nom\,\K\,\tau$: cupón por periodo; $\mathit{CF}_i$: flujo $i$ (moneda del bono).
  \item $\mathrm{AI}$: cupón corrido; $t_{\mathrm{prev}}\le0$: tiempo del último cupón ya pagado respecto de la
        fecha de valuación.
  \item $V$ (\PV): precio sucio, $V_{\mathrm{lim}}=V-\mathrm{AI}$: precio limpio.
  \item $y$: rendimiento con composición continua (tasa anual); $D$: duración (años); $\mathrm{Cvx}$:
        convexidad (años$^{2}$).
  \item En el envejecimiento: $e$ días transcurridos desde la fecha de inicio; $d_i$ día del pago $i$; $i_1$
        índice del primer pago pendiente.
\end{itemize}

% -------------------------------------------------------------------------------------------------
\section{Modelo y fórmulas}\label{sec:cap06-modelo}

\subsection{Flujos y precio sucio}\label{sec:cap06-flujos}

Con $m$ periodos y frecuencia $f$ el constructor regular de \codigo{include/qf/instruments/bond.hpp}{Bond}
fija los pagos en $t_i=i/f$ ($i=1,\dots,m$) y los flujos
\begin{equation}\label{eq:cap06-flujos}
  C=\Nom\,\K\,\tau,\qquad \mathit{CF}_i=C\ (i<m),\qquad \mathit{CF}_m=C+\Nom,
\end{equation}
con $\tau=\texttt{periodFraction}(f,\mathrm{dcc})$ de \codigo{src/math/daycount.cpp}{periodFraction}: $1/f$ bajo
ACT/365 (la base por omisión, y la única que usa el desk), $365/(360f)$ bajo ACT/360. No hay conteo de días
reales: el cupón es el mismo en cada periodo (cálculo en \codigo{src/instruments/bond.cpp}{Bond::cashflows}).
El precio sucio es el valor presente de los flujos pendientes en la curva de descuento,
\begin{equation}\label{eq:cap06-precio}
  V=\sum_{i}\mathit{CF}_i\,\DF{c}{t_i},
\end{equation}
calculado por \codigo{src/instruments/bond.cpp}{Bond::price}; \codigo{src/instruments/bond.cpp}{Bond::calculatePV}
lo toma con la curva predeterminada del entorno de mercado. Los flujos que ve el desk son los de
\codigo{src/instruments/bond.cpp}{Bond::cashFlowSchedule}, pares (tiempo en años, monto) en los que el último
monto incluye el principal.

\subsection{Calendario explícito y cupón corrido}\label{sec:cap06-corrido}

El segundo constructor recibe los tiempos de pago pendientes $t_1<\dots<t_n$ (finitos, estrictamente
crecientes y positivos) y el tiempo $t_{\mathrm{prev}}\le0$ del último cupón pagado. El cupón $C$ sigue
siendo el de la ecuación~\eqref{eq:cap06-flujos} (la frecuencia solo fija su tamaño) y el último pago incluye el
principal. El cupón corrido es la parte transcurrida del periodo en curso, lineal en el tiempo, y el precio
limpio es el sucio menos el corrido:
\begin{equation}\label{eq:cap06-corrido}
  \mathrm{AI}=C\,\frac{-t_{\mathrm{prev}}}{t_1-t_{\mathrm{prev}}}\quad(\mathrm{AI}=0\text{ si }t_{\mathrm{prev}}=0),
  \qquad V_{\mathrm{lim}}=V-\mathrm{AI}.
\end{equation}
Se implementa en \codigo{src/instruments/bond.cpp}{Bond::accruedInterest} y
\codigo{src/instruments/bond.cpp}{Bond::cleanPrice}; el constructor regular siempre tiene $\mathrm{AI}=0$ y
su precio limpio coincide con el sucio. Con $t_{\mathrm{prev}}=0$ y pagos en $i/f$ el segundo constructor
reproduce al primero.

\subsection{Rendimiento, duración y convexidad}\label{sec:cap06-rendimiento}

El rendimiento $y$ que devuelve \codigo{src/instruments/bond.cpp}{Bond::yield} resuelve, para un precio $P$
dado,
\begin{equation}\label{eq:cap06-rendimiento}
  \sum_{i}\mathit{CF}_i\,e^{-y\,t_i}=P,
\end{equation}
con composición continua y sin descontar con ninguna curva: el precio $P$ se compara con la suma de todos los
flujos pendientes, así que es un precio sucio. El código encierra la raíz en $[-0.30,\,0.10]$, desplaza el extremo
superior de $0.10$ en $0.10$ mientras la función siga positiva (hasta $50$) y el inferior de $0.10$ en $0.10$
mientras siga negativa (hasta $-5$), lanza error si no logra un cambio de signo y aplica el método de Brent
(\codigo{include/qf/math/rootfinding.hpp}{brent}). El binding Python es
\codigo{src/python_bindings/qfpy.cpp}{yield}, pero el desk no lo usa.

La duración y la convexidad se miden contra la curva de descuento, no contra $y$:
\begin{equation}\label{eq:cap06-duracion}
  D=\frac{1}{V}\sum_{i}t_i\,\mathit{CF}_i\,\DF{c}{t_i},\qquad
  \mathrm{Cvx}=\frac{1}{V}\sum_{i}t_i^{2}\,\mathit{CF}_i\,\DF{c}{t_i},
\end{equation}
en \codigo{src/instruments/bond.cpp}{Bond::duration} y \codigo{src/instruments/bond.cpp}{Bond::convexity}.
Son el promedio ponderado por valor presente del tiempo de los flujos (duración de Macaulay) y su
análogo de segundo orden; también son $-\frac1V\partial V/\partial s$ y $\frac1V\partial^2V/\partial s^2$
para un desplazamiento paralelo $s$ de la tasa cero continua de la curva. El comentario del código afirma que
con composición continua la duración modificada coincide con la de Macaulay.

\begin{diferencia}
  En el mercado los bonos se cotizan por un rendimiento periódico (por ejemplo, semestral equivalente a
  bono) y la duración modificada es $D_{\mathrm{Mac}}/(1+y/f)$ \citep{hull2018,tuckman2012}. Aquí el
  rendimiento es continuo, $y=f\ln\!\bigl(1+y_{\mathrm{per}}/f\bigr)$ o, al revés,
  $y_{\mathrm{per}}=f\bigl(e^{y/f}-1\bigr)$, y la «duración» que el desk reporta (pantalla de pricers) es la de
  Macaulay, que bajo composición continua es la modificada; con rendimiento periódico no lo sería. El
  rendimiento tampoco se usa para valuar: el valor viene siempre de la curva.
\end{diferencia}

\begin{diferencia}
  El cupón es $\Nom\K/f$ en todos los periodos, sin contar los días reales de cada uno (la base ACT/365 por
  omisión da $\tau=1/f$). La convención M-Bono de 182 días queda fuera (Metodología, «Límites», «Bono»). La Metodología, en «Convención de días», limita ahora el devengo ACT/360 a los swaps y a las opciones de tasa y
  declara este cupón (fila «Cupón» del bono): hallazgo H13 de \texttt{HALLAZGOS.md}, resuelto.
\end{diferencia}

\subsection{Tabla de convenciones}

\begin{table}[ht]
\centering
\begin{tabular}{>{\raggedright\arraybackslash}p{0.28\textwidth}>{\raggedright\arraybackslash}p{0.58\textwidth}}
\toprule
Aspecto & Qué hace el código \\
\midrule
Cupón & $\Nom\K\tau$ por periodo, igual en todos; $\tau=1/f$ (ACT/365 por omisión) \\
Tiempo de pago & $t_i=i/f$ años (bono nuevo); $i/f-e/365$ años (envejecido) \\
Descuento & $\DF{c}{t_i}$ de la curva de la operación o de la del colateral (sección~\ref{sec:cap06-desk}) \\
Precio del producto & sucio (en el bono nuevo sucio y limpio coinciden) \\
Rendimiento & composición continua, sobre el precio sucio; no lo usa el desk \\
Duración y convexidad & de la curva de descuento, ponderadas por valor presente \\
Corrido & lineal en tiempo exacto; 0 si el último cupón cae hoy \\
\bottomrule
\end{tabular}
\caption{Convenciones del bono.}\label{tab:cap06-convenciones}
\end{table}

% -------------------------------------------------------------------------------------------------
\section{Implementación}\label{sec:cap06-implementacion}

\subsection{Del producto \texttt{bond} al instrumento}\label{sec:cap06-desk}

Los pricers de tasas se registran en \codigo{desk/api/pricing/rates.py}{PRODUCTS}. Para el bono:
\begin{itemize}
  \item \codigo{desk/api/pricing/rates.py}{_bond_validate} normaliza los parámetros: curva (por omisión
        \texttt{TIIE-F}), \texttt{face}~$>0$, \texttt{coupon\_rate} en $[0,\text{MAX\_RATE}]$, \texttt{periods}
        entero $\ge1$, \texttt{frequency} por omisión $2$ y rechaza duraciones $m/f$ de más de 1200 años.
  \item \codigo{desk/api/pricing/rates.py}{_bond_instrument} construye
        \texttt{qfpy.Bond(face, coupon\_rate, periods, frequency)} (constructor regular, ACT/365).
  \item La curva de descuento es la de la operación, salvo que sea \texttt{TIIE-F} con colateral en USD y
        exista la curva \texttt{MXN-USD}, que entonces descuenta (capítulo~\ref{cap:curvas}); \texttt{SOFR}
        siempre descuenta con SOFR. \codigo{desk/api/pricing/rates.py}{_bond_pv} evalúa
        $V$ con el entorno cuya curva predeterminada es la de descuento.
  \item \codigo{desk/api/pricing/rates.py}{_bond_price} devuelve, además de $V$, la tabla de flujos (días
        $=\operatorname{round}(365\,t_i)$, fecha, factor de descuento calculado con $t_i$ exacto, y valor
        presente), el DV01 ($V(\text{base})-V(+1~\pb)$ con un choque paralelo a la curva de la operación) y
        $D$ y $\mathrm{Cvx}$ de la ecuación~\eqref{eq:cap06-duracion}.
  \item \codigo{desk/api/pricing/rates.py}{_bond_describe} informa el vencimiento como la fecha de valuación más
        $\operatorname{round}(365\,t_m)$ días.
\end{itemize}
El PV está en la moneda de la curva (\MXN\ para \texttt{TIIE-F}, \USD\ para \texttt{SOFR}); el pricer de
producto lo convierte a \MXN\ con el spot del snapshot.

\subsection{Bono envejecido}\label{sec:cap06-envejecimiento}

Para valuar un bono en una fecha posterior a su inicio, \codigo{desk/api/valuation.py}{value_trade} llama a
\codigo{desk/api/aging.py}{bond_pv_mxn}, que a su vez arma el instrumento con calendario explícito en
\codigo{desk/api/aging.py}{aged_bond}. Sean $e$ los días entre la fecha de inicio y la de valuación y
\begin{equation}\label{eq:cap06-dias-pago}
  d_i=\operatorname{round}\!\bigl(365\,i/f\bigr),\qquad i=1,\dots,m,
\end{equation}
el día del pago $i$ (\codigo{desk/api/aging.py}{_pay_day}; la misma fecha que el vencimiento de la operación).
Quedan por pagar los pagos con $d_i>e$ (\codigo{desk/api/aging.py}{_remaining_payments}); un cupón con fecha
igual a la de valuación ya está liquidado y no entra. Sus tiempos son \emph{exactos}, no redondeados a días:
\begin{equation}\label{eq:cap06-tiempos-envejecido}
  t_k=\frac{i_k}{f}-\frac{e}{365},\qquad
  t_{\mathrm{prev}}=\min\Bigl\{0,\ \frac{i_1-1}{f}-\frac{e}{365}\Bigr\},
\end{equation}
donde $i_1<i_2<\dots$ son los índices de los pagos pendientes. Con la base ACT/365 del cupón,
la ecuación~\eqref{eq:cap06-corrido} se reduce a $\mathrm{AI}=\Nom\K\,(-t_{\mathrm{prev}})$: el cupón corrido es
la tasa anual por el tiempo transcurrido desde el último cupón (años de $i/f$ y de días$/365$), y es cero en el
día en que cae un cupón. Si el bono aún no empieza ($e<0$), los tiempos se alargan en $|e|/365$ y
$\mathrm{AI}=0$. El recorte a $0$ en $t_{\mathrm{prev}}$ evita que el redondeo del día deje el cupón anterior
unas horas después de la fecha de valuación, que C++ rechazaría. Si no queda ningún pago, \texttt{aged\_bond}
lanza \texttt{ValueError}.

Las funciones de salida son \codigo{desk/api/aging.py}{bond_pv_mxn} (precio sucio, convertido con el spot si la
curva es SOFR), \codigo{desk/api/aging.py}{bond_quote} (sucio, limpio y corrido en la moneda de la curva, con
$\text{sucio}-\text{limpio}=\text{corrido}$) y \codigo{desk/api/aging.py}{projected_bond} (flujos por cobrar
con factor de descuento y valor presente). La liquidación de los cupones y de la amortización, con fechas
$\text{inicio}+d_i$ y cupón $C$ tomado del constructor de C++
(\codigo{desk/api/aging.py}{_coupon_amount}), vive en \codigo{desk/api/aging.py}{bond_settled}
(capítulo~\ref{cap:ciclo-vida}).

Dos consecuencias que la Metodología declara y que conviene tener presentes. Primero, el calendario de
\emph{valuación} usa tiempos exactos $i/f$ mientras la \emph{liquidación} cae en el día entero $d_i$; con
$f=2$, $d_1=\operatorname{round}(182.5)=182$, de modo que el día siguiente al cupón el corrido es de medio día
y no de un día completo. Segundo, \texttt{round} de Python redondea los empates al entero par:
$d_1=182$ pero $d_3=\operatorname{round}(547.5)=548$, así que el primer cupón se paga medio día antes y el tercero
medio día después de la fecha exacta.

% -------------------------------------------------------------------------------------------------
\section{Ejemplo numérico verificable}\label{sec:cap06-ejemplo}

Los valores se reproducen con \codigo{docs/tex/verificacion/cap06_bonos.py}{bono_regular} y
\codigo{docs/tex/verificacion/cap06_bonos.py}{bono_envejecido}; el script llama a \texttt{qfpy.Bond} y a
\texttt{desk.api.aging} y comprueba cada valor con la tolerancia indicada.

\subsection{Bono de 2 años, cupón 8\% semestral, curva plana a 6\%}

$\Nom=10^{6}$, $\K=8\%$, $f=2$, $m=4$, curva plana a $6\%$ continuo. Por la
ecuación~\eqref{eq:cap06-flujos}, $C=40\,000$ y los flujos son $40\,000$ en $t=0.5,\,1,\,1.5$ y $1\,040\,000$ en
$t=2$. Con $\DF{c}{t}=e^{-0.06t}$, es decir $0.9704455$, $0.9417645$, $0.9139312$ y $0.8869204$, la
ecuación~\eqref{eq:cap06-precio} da
\[
  V = 40\,000\,(0.9704455+0.9417645+0.9139312)+1\,040\,000\cdot0.8869204
  = \verif{cap06-precio-sucio}{1035442.9043}\ \MXN,
\]
con tolerancia absoluta $10^{-4}$. La duración y la convexidad de la ecuación~\eqref{eq:cap06-duracion} son
\[
  D=\verif{cap06-duracion}{1.88973227}\ \text{años},\qquad
  \mathrm{Cvx}=\verif{cap06-convexidad}{3.68848716}\ \text{años}^2
\]
(tolerancia absoluta $10^{-8}$).

\subsection{Rendimiento a la par}

Si el precio fuera el nominal ($P=10^{6}$) el rendimiento de la ecuación~\eqref{eq:cap06-rendimiento} es
\[
  y=\verif{cap06-rendimiento-par}{0.0784414263}\quad(\text{tolerancia absoluta }10^{-9}),
\]
y su equivalente semestral $y_{\mathrm{per}}=2\bigl(e^{y/2}-1\bigr)=\verif{cap06-rendimiento-semestral}{0.0800000000}$
(tolerancia absoluta $10^{-9}$): a la par el rendimiento periódico es el cupón, mientras que el continuo,
$7.844\%$, es menor. Ambos números describen el mismo bono; solo difiere la composición (recuadro
«Diferencia» de la sección~\ref{sec:cap06-rendimiento}).

\subsection{Bono envejecido 90 días}

Mismo bono con inicio el 2026-01-02 y fecha de valuación 90 días después ($e=90$), en la curva plana a $6\%$.
Los cuatro pagos están pendientes ($d_i=182,\,365,\,548,\,730>90$ con $f=2$), con tiempos
$t_k=i/2-90/365=0.2534247,\,0.7534247,\,1.2534247,\,1.7534247$ y $t_{\mathrm{prev}}=0-90/365=-0.2465753$. El
corrido de la ecuación~\eqref{eq:cap06-corrido} es $C\cdot0.2465753/0.5=\Nom\K\cdot90/365$:
\[
  \mathrm{AI}=\verif{cap06-corrido}{19726.0274}\ \MXN
\]
(tolerancia absoluta $10^{-4}$). Con curva plana y ningún cupón pagado todavía, el sucio es el de la fecha de
pacto capitalizado por 90 días,
$V=1\,035\,442.9043\,e^{0.06\cdot90/365}=\verif{cap06-precio-sucio-envejecido}{1050875.6643}$ \MXN, y el limpio es
$V-\mathrm{AI}=\verif{cap06-precio-limpio}{1031149.6369}$ \MXN (ambos con tolerancia absoluta $10^{-4}$; el script
comprueba la capitalización de forma independiente).

\subsection{Medio día de corrido}

En $e=183$ el pago $i=1$ ($d_1=182$) ya está liquidado, quedan tres pagos y $t_{\mathrm{prev}}=0.5-183/365=
-0.0013699$ (el cupón exacto cae a los $182.5$ días, medio día antes de la valuación), con
\[
  \mathrm{AI}=\Nom\K\cdot\frac{0.5}{365}=\verif{cap06-corrido-dia-183}{109.5890}\ \MXN
\]
(tolerancia absoluta $10^{-4}$): medio día de corrido, no un día completo.

% -------------------------------------------------------------------------------------------------
\section{Validación}\label{sec:cap06-validacion}

\begin{itemize}
  \item Pruebas de C++: el calendario regular reproduce al constructor por periodos, el corrido es la parte
        transcurrida del cupón, limpio más corrido es el sucio, el sucio se mueve por el carry diario entre
        cupones y el limpio no da saltos al pagar un cupón; usan la función auxiliar
        \codigo{tests/test_bond_schedule.cpp}{agedBond}, que replica el envejecimiento del desk. Otras pruebas
        cubren el rendimiento (ida y vuelta, y a la par), la duración, la convexidad y los flujos; sus curvas de
        prueba son \codigo{tests/test_bond_cashflows.cpp}{upwardCurve} y
        \codigo{tests/test_instruments.cpp}{flatCurve}.
  \item Bindings:
        \codigo{tests/test_aged_bond_cap_py.py}{AgedBondTests.test_accrued_interest_and_clean_price}.
  \item Backend: el PV envejecido el día del pacto coincide con el precio y no hay escalones al rodar día a día
        (\codigo{desk/api/tests/test_aging_exact.py}{test_bond_rolled_day_by_day_has_no_step_and_dirty_clean_accrued_add_up}),
        el corrido crece linealmente entre cupones
        (\codigo{desk/api/tests/test_aging_exact.py}{test_bond_accrued_interest_grows_linearly_between_coupons}),
        las fechas de pago coinciden con el vencimiento para cualquier frecuencia
        (\codigo{desk/api/tests/test_aging_exact.py}{test_bond_pay_days_match_the_maturity_date_for_any_frequency})
        y la forma de la respuesta del pricer y el bono cupón cero
        (\codigo{desk/api/tests/test_price_rates.py}{test_bond_response_keys_and_cashflows},
        \codigo{desk/api/tests/test_price_rates.py}{test_zero_coupon_bond_prices}).
  \item Contra QuantLib (capítulo~\ref{cap:validacion}): \codigo{desk/validation/ql_rates.py}{price_bond} valúa
        el mismo bono con \texttt{FixedRateBond} (ACT/365) y \texttt{DiscountingBondEngine} sobre la curva del
        desk, con pagos en $\operatorname{round}(365\,i/f)$ días enteros. Los 8 casos caen dentro de la
        tolerancia de $10^{-6}$ del nocional; la mayor diferencia relativa es $5.9\times10^{-8}$ (caso
        \codigo{desk/validation/report.json}{bond-08}, con cupones cada 182.5 días), del orden del efecto de
        redondear las fechas de pago a días enteros.
\end{itemize}

% -------------------------------------------------------------------------------------------------
\section{Límites y supuestos}\label{sec:cap06-limites}

\begin{limites}
  \begin{itemize}
    \item Solo bonos de tasa fija con cupones iguales y amortización total al vencimiento («Periodos
          uniformes», Metodología, «Límites»); la convención de los bonos mexicanos de 182 días queda fuera.
    \item El cupón no usa una base de días: $\Nom\K/f$; el corrido es lineal en tiempo exacto (días$/365$). No hay
          ajuste por días inhábiles.
    \item El valor sale de descontar con la curva swap elegida (\texttt{TIIE-F}, \texttt{SOFR}, o la de
          colateral): no hay curva ni diferencial de crédito del emisor, ni se valúa por rendimiento.
    \item El rendimiento es continuo y la duración es la de Macaulay (la modificada bajo composición continua);
          no son las medidas periódicas del mercado.
    \item En el envejecimiento la valuación usa tiempos exactos $i/f$ y la liquidación usa el día entero más
          cercano, con empates al par; el corrido el día siguiente a un cupón semestral es de medio día
          (Metodología, «Bono y cap/floor envejecidos»).
    \item El PV del bono en el portafolio es el precio \emph{sucio}; el limpio y el corrido salen de
          \codigo{desk/api/aging.py}{bond_quote}.
    \item Los parámetros del producto están acotados (\texttt{MAX\_PERIODS}, \texttt{MAX\_BOND\_YEARS}); un bono
          ya vencido no tiene flujos por valuar.
  \end{itemize}
\end{limites}

% -------------------------------------------------------------------------------------------------
\section{Referencias}\label{sec:cap06-referencias}

Las definiciones de precio sucio y limpio, cupón corrido, rendimiento y duración siguen a
\citet{tuckman2012} y a \citet{hull2018}; las diferencias con la práctica de mercado de la
sección~\ref{sec:cap06-rendimiento} se refieren a esas dos obras.
